% OpenStax Calculus Volume 2 -- Section 2.3 Exercises: Volumes of Revolution: Cylindrical Shells
% Exercises reproduced from OpenStax, licensed under CC BY 4.0.
% Source: https://openstax.org/books/calculus-volume-2/pages/2-3-volumes-of-revolution-cylindrical-shells
\documentclass{exercises}
\title{OpenStax Calculus Volume 2}
\subtitle{Section 2.3 Exercises: Volumes of Revolution: Cylindrical Shells}
\sourceurl{https://openstax.org/books/calculus-volume-2/pages/2-3-volumes-of-revolution-cylindrical-shells}
\author{}
\date{}

\begin{document}
\maketitle


For the following exercises, find the volume generated when the region between the two curves is rotated around the given axis. Use both the shell method and the washer method. Use technology to graph the functions and draw a typical slice by hand.

\begin{enumerate}[start=1]
\item \techrequired Bounded by the curves $y=3x,x=0,\,\text{and}\,y=3$ rotated around the $y\text{-axis}$.
\item \techrequired Bounded by the curves $y=3x,y=0,\,\text{and}\,x=3$ rotated around the $y\text{-axis}$.
\item \techrequired Bounded by the curves $y=3x,y=0,\,\text{and}\,y=3$ rotated around the $x\text{-axis}$.
\item \techrequired Bounded by the curves $y=3x,y=0,\,\text{and}\,x=3$ rotated around the $x\text{-axis}$.
\item \techrequired Bounded by the curves $y=2x^{3},y=0,\,\text{and}\,x=2$ rotated around the $y\text{-axis}$.
\item \techrequired Bounded by the curves $y=2x^{3},y=0,\,\text{and}\,x=2$ rotated around the $x\text{-axis}$.
\end{enumerate}

For the following exercises, use shells to find the volumes of the given solids. Note that the rotated regions lie between the curve and the $x\text{-axis}$ and are rotated around the $y\text{-axis}$.

\begin{enumerate}[resume]
\item $y=1-x^{2},x=0,\,\text{and}\,x=1$
\item $y=5x^{3},x=0,\,\text{and}\,x=1$
\item $y=\frac{1}{x},x=1,\,\text{and}\,x=100$
\item $y=\sqrt{1-x^{2}},x=0,\,\text{and}\,x=1$
\item $y=\frac{1}{1+x^{2}},x=0,\,\text{and}\,x=3$
\item $y=\sin(x^{2}),x=0,\,\text{and}\,x=\sqrt{\pi}$
\item $y=\frac{1}{\sqrt{1-x^{2}}},x=0,\,\text{and}\,x=\frac{1}{2}$
\item $y=\sqrt{x},x=0,\,\text{and}\,x=1$
\item $y=(1+x^{2})^{3},x=0,\,\text{and}\,x=1$
\item $y=5x^{3}-2x^{4},x=0,\,\text{and}\,x=2$
\end{enumerate}

For the following exercises, use shells to find the volume generated by rotating the regions between the given curve and $y=0$ around the $x\text{-axis}$.

\begin{enumerate}[resume]
\item $y=\sqrt{1-x^{2}},x=0,\,x=1$ and the $x\text{-axis}$
\item $y=x^{2},x=0,\,x=2$ and the $x\text{-axis}$
\item $y=\frac{x^{3}}{2},\,x=0,\,x=2$, and the $x\text{-axis}$
\item $y=\frac{2}{x^{2}},\,x=1,\,x=2$, and the $x\text{-axis}$
\item $x=\frac{1}{1+y^{2}},y=4$
\item $x=\frac{1+y^{2}}{y},y=1,\,y=4$, and the $y\text{-axis}$
\item $x=\sqrt{4-y^{2}},x=0,y=0$
\item $x=y^{3}-2y^{2},\,x=0,\,x=9$
\item $x=\sqrt{y}+1,\,x=1,\,x=3$, and the $x\text{-axis}$
\item $x=\sqrt[3]{27y}\,\text{and}\,x=\frac{3y}{4}$
\end{enumerate}

For the following exercises, find the volume generated when the region between the curves is rotated around the given axis.

\begin{enumerate}[resume]
\item $y=3-x,y=0,x=0,\,\text{and}\,x=2$ rotated around the $y\text{-axis}$.
\item $y=x^{3},x=0,\,\text{and}\,y=8$ rotated around the $y\text{-axis}$.
\item $y=x^{2},y=x$, rotated around the $y\text{-axis}$.
\item $y=\sqrt{x},y=0,\,\text{and}\,x=1$ rotated around the line $x=2$.
\item $y=\frac{1}{4-x},x=1,\,x=2\,\text{and}\,y=0$ rotated around the line $x=4$.
\item $y=\sqrt{x}\,\text{and}\,y=x^{2}$ rotated around the $y\text{-axis}$.
\item $y=\sqrt{x}\,\text{and}\,y=x^{2}$ rotated around the line $x=2$.
\item $x=y^{3},x=\frac{1}{y},x=1,\,\text{and}\,x=2$ rotated around the $x\text{-axis}$.
\item $x=y^{2}\,\text{and}\,y=x$ rotated around the line $y=2$.
\item \techrequired Left of $x=\sin(\pi y)$, right of $y=x$, around the $y\text{-axis}$.
\end{enumerate}

For the following exercises, use technology to graph the region. Determine which method you think would be easiest to use to calculate the volume generated when the function is rotated around the specified axis. Then, use your chosen method to find the volume.

\begin{enumerate}[resume]
\item \techrequired $y=x^{2}$ and $y=4x$ rotated around the $y\text{-axis}$.
\item \techrequired $y=\cos(\pi x),y=\sin(\pi x),x=\frac{1}{4},\,\text{and}\,x=\frac{5}{4}$ rotated around the $y\text{-axis}$. This exercise requires advanced technique. You may use technology to perform the integration.
\item \techrequired $y=x^{2}-2x,x=2,\,\text{and}\,x=4$ rotated around the $y\text{-axis}$.
\item \techrequired $y=x^{2}-2x,x=2,\,\text{and}\,x=4$ rotated around the $x\text{-axis}$.
\item \techrequired $y=3x^{3}-2,y=x,\,\text{and}\,x=2$ rotated around the $x\text{-axis}$.
\item \techrequired $y=3x^{3}-2,y=x,\,\text{and}\,x=2$ rotated around the $y\text{-axis}$.
\item \techrequired $x=\sin(\pi y^{2})$ and $x=\sqrt{2}y$ rotated around the $x\text{-axis}$.
\item \techrequired $x=y^{2},x=y^{2}-2y+1,\,\text{and}\,x=2$ rotated around the $y\text{-axis}$.
\end{enumerate}

For the following exercises, use the method of shells to approximate the volumes of some common objects, which are pictured in accompanying figures.

\begin{enumerate}[resume]
\item Use the method of shells to find the volume of a sphere of radius $r$. \\ \textit{[Figure: This figure has two images. The first is a circle with radius $r$. The second is a basketball.]}
\item Use the method of shells to find the volume of a cone with radius $r$ and height $h$. \\ \textit{[Figure: This figure has two images. The first is an upside-down cone with radius $r$ and height $h$. The second is an ice cream cone.]}
\item Use the method of shells to find the volume of an ellipsoid $(x^{2}/a^{2})+(y^{2}/b^{2})=1$ rotated around the $x\text{-axis}$. \\ \textit{[Figure: This figure has two images. The first is an ellipse with $a$ the horizontal distance from the center to the edge and $b$ the vertical distance from the center to the top edge. The second is a watermelon.]}
\item Use the method of shells to find the volume of a cylinder with radius $r$ and height $h$. \\ \textit{[Figure: This figure has two images. The first is a cylinder with radius $r$ and height $h$. The second is a cylindrical candle.]}
\item Use the method of shells to find the volume of the donut created when the circle $x^{2}+y^{2}=4$ is rotated around the line $x=4$. \\ \textit{[Figure: This figure has two images. The first has two ellipses, one inside of the other. The radius of the path between them is 2 units. The second is a doughnut.]}
\item Consider the region enclosed by the graphs of $y=f(x),y=1+f(x),x=0$, and $x=a>0$. What is the volume of the solid generated when this region is rotated around the $y\text{-axis}$? Assume that the function is defined over the interval $[0,a]$.
\item Consider the function $y=f(x)$, which decreases from $f(0)=b$ to $f(1)=0$. Set up the integrals for determining the volume, using both the shell method and the disk method, of the solid generated when this region, with $x=0$ and $y=0$, is rotated around the $y\text{-axis}$. Prove that both methods approximate the same volume. Which method is easier to apply? (Hint: Since $f(x)$ is one-to-one, there exists an inverse $f^{-1}(y)$.)
\end{enumerate}

\end{document}
